\scanpage{37}
in other words $X$ and $\bar X$ have the same sheaf theory:
For all \inserted{(most?)} practical purposes, $X$ can be replaced with $\widetilde X$.

\deleted{\unsure{it}} We have a map

\[
\operatorname{Hom}_{\mathrm{Sp}}^{\unsure{\prime\prime}}(X,Y)\longrightarrow
\operatorname{Hom}_{\mathrm{top}}(\operatorname{Top}(X),\operatorname{Top}(Y))
\qquad(\text{isom classes})
\]

it is bijective if $Y$ is sober. In fact, the category $\underline{\mathrm{Hom}}_{\mathrm{top}}(\operatorname{Top}(X),\operatorname{Top}(Y))$
is \deleted{discrete}\inserted{rigid}, associated to an order relation on $\operatorname{Hom}_{\mathrm{top}}(\operatorname{Top}(X),\operatorname{Top}(Y))$
(specialisation) !

\uline{The notion of a topos}. \inserted{(gener\unsure{.?|al.?|.}) ``topology'' $=$ study of topoi,}
Why it is a natural and useful generalisation of the notion of topological space?

a) It is the natural setting for most ``algebraic'' (and non algebraic) structures.

b) As such, it allows to construct classifying topoi for such structures (possibly all): algebra = topology !

c) It is stable under a variety of operations that in many cases do not make sense in the category of topological spaces:
$*$ generalized notions of passage to quotient (by wild equivalence relations, which would give pathological quotients within top. spaces)

$*$ clutching together of pieces à la Artin

d) They arise in various questions \inserted{of geometry} as modular topoi

e) They provide the most convenient set-up for
\inserted{all topics involving systematically} \deleted{\unsure{this}} passing to quotients.

f) They give the ultimate outcome for \inserted{the intuitive notion of} localization.

g) They are technical musts for various cohomology theories
in algebraic geometry.

\uline{Examples of topoi}

\[
\left\{\begin{aligned}
&\operatorname{Top}(X),\quad \operatorname{Top}(X,G)
 \left\{\begin{aligned}&\operatorname{Top}(\unsure{X|X'})\\&B_G\quad H^{\unsure{q|*}}(G,-)\end{aligned}\right.\\
&\widehat C\quad\text{for instance }B_G\ (G\text{ discrete})
\end{aligned}\right.
\]
% archive sheet number: 36
